% ======================================================================
%  TIER 2 KARTI: Skill (Agent Skills)
% ======================================================================

\card{Skill (Agent Skills)}{ihtiyaç anında yüklenen paketli yetkinlik}
\cl{Problem} Bir agent onlarca işi yapabilmeli --- ama hepsinin talimatını,
script'ini ve örneğini aynı anda system prompt'a veya araç listesine gömersen
pencere şişer, tekrar kullanılamaz ve bakımı kâbusa döner.
\cl{Tanım} \term{Skill}, bir agent'ın \emph{ihtiyaç anında} yüklediği,
paketlenmiş ve yeniden kullanılabilir bir yetkinliktir: talimat + (gerekirse)
script ve kaynaklar. İlgili olmayan skill context'e hiç girmez (\en{progressive disclosure}).
\cl{Model} Uzmanın masasında her konunun el kitabı durmaz; o işe başlarken ilgili
kitapçığı açar. Skill, ihtiyaç anında açılan o el kitabıdır.
\cl{Ne zaman} Tekrar eden, paketlenebilir uzmanlıkta (``PDF doldur'', ``marka
tonuyla yaz''); çok sayıda yeteneği bir agent'a, context'i şişirmeden vermek için.
\cl{Skill vs Tool vs MCP} Tek fonksiyon → \term{Tool}; çapraz-uygulama bağlantı →
\term{MCP}; talimat + kaynak içeren paketli yetkinlik → \textbf{Skill}.
\cl{Tuzak} Her şeyi skill yapmak; basit tek-fonksiyonu skill'e sarmak; yüklenince
devasa talimatla context'i şişirmek.
\cl{İlişkili} \term{Tool Calling} · \term{MCP} · \term{Context Engineering} · \term{Agent} \,{\scriptsize\color{softgray}(İlke \#3, \#6)}
